%% ======================================================================
%% C3 Graduated Filtering — Contingency Text Versions
%% Pre-written for exp-006 results (ETA 2026-05-15).
%% Replace XX.X ± X.X placeholders with actual numbers once available.
%% ======================================================================


%% ======================================================================
%% VERSION A — Graduated filtering significantly outperforms binary
%% ======================================================================

%% --- For conclusion.tex (replace the C3 TODO block, lines 48–53) ---

% === Contribution 3: Graduated filtering (C3) ===
% Source: exp-006
The $\Kempirical = 2$ boundary implies a natural trichotomy---human
(depth~0), first-generation synthetic (depth~1), and deeper synthetic
(depth~$\geq 2$)---that can be exploited for depth-aware data curation.
Graduated filtering, which assigns sample weights $w = \exp(-\alpha
\cdot \hat{d})$ based on the OBD-predicted expected depth $\hat{d}$,
outperforms both binary filtering strategies on downstream benchmarks:
at the optimal $\alpha^* = $ XX.X, LoRA-tuned Qwen2.5-1.5B-Instruct
achieves XX.X $\pm$ X.X on MMLU and XX.X $\pm$ X.X on HellaSwag,
compared to XX.X $\pm$ X.X / XX.X $\pm$ X.X for binary-ours and
XX.X $\pm$ X.X / XX.X $\pm$ X.X for binary-Drayson
(Table~\ref{tab:filtering_main}).
The gain is attributable to the retention of depth-1 content, which
preserves substantial original signal while only gently down-weighting
first-generation text; binary strategies discard this stratum entirely,
sacrificing data volume without commensurate quality improvement.
This validates the practical utility of the three-tier taxonomy derived
from $\Kempirical = 2$: depth-aware curation that distinguishes shallow
from deep synthetic content yields measurably better training corpora
than the binary human-versus-synthetic partition.

%% --- For discussion.tex (new subsection or append to §Broader Impact) ---

\paragraph{Graduated versus binary filtering.}
The graduated filtering results (Section~\ref{sec:graduated-filtering})
provide direct evidence that the $\Kempirical = 2$ trichotomy has
actionable value for data curation.
The optimal decay parameter $\alpha^* = $ XX.X falls in the moderate
range, neither so low as to approximate no-filter ($\alpha \to 0$) nor
so high as to degenerate to hard binary filtering ($\alpha \to \infty$).
This intermediate optimum reflects the core insight: depth-1 text, while
detectably synthetic, retains enough of the original distribution's
diversity---as evidenced by elevated surprisal variance and vocabulary
richness relative to depth~$\geq 2$ (Section~\ref{sec:feature_degradation})---to
benefit downstream training when appropriately down-weighted rather than
discarded.
The MMLU improvement of $+$XX.X\,pp over binary-ours and
$+$XX.X\,pp over binary-Drayson, together with consistent gains on
HellaSwag ($+$XX.X\,pp and $+$XX.X\,pp, respectively), demonstrate
that the signal preserved in the depth-1 stratum translates to
measurable downstream quality across both knowledge-intensive and
commonsense reasoning tasks.
Graduated filtering also exhibits lower cross-seed variance than binary
strategies (XX.X versus XX.X on MMLU), suggesting that the ordinal
weighting scheme provides a more stable curation signal than hard
classification boundaries.


%% ======================================================================
%% VERSION B — Graduated filtering NOT significantly better than binary
%% ======================================================================

%% --- For conclusion.tex (replace the C3 TODO block, lines 48–53) ---

% === Contribution 3: Graduated filtering (C3) ===
% Source: exp-006
The $\Kempirical = 2$ boundary implies a natural trichotomy---human
(depth~0), first-generation synthetic (depth~1), and deeper synthetic
(depth~$\geq 2$)---that enables testing whether ordinal depth information
improves data curation beyond a simple binary filter.
We evaluate graduated filtering, which assigns sample weights
$w = \exp(-\alpha \cdot \hat{d})$ based on predicted expected depth, against
binary filtering baselines (Table~\ref{tab:filtering_main}).
At the best $\alpha$, graduated filtering achieves XX.X $\pm$ X.X on
MMLU and XX.X $\pm$ X.X on HellaSwag, comparable to binary-ours
(XX.X $\pm$ X.X / XX.X $\pm$ X.X) and binary-Drayson
(XX.X $\pm$ X.X / XX.X $\pm$ X.X); the differences are not
statistically significant across 5 seeds ($p >$ XX.XX).
This result is not an experimental limitation but a direct consequence
of $\Kempirical$'s status as an operationally sufficient boundary:
binary filtering at the $\Kempirical$ threshold already captures the
actionable depth signal, and graduated weighting cannot improve on it
because the residual ordinal information beyond the binary partition
is insufficient to meaningfully differentiate training quality.
Binary filtering at $\Kempirical$ is the simplest optimal curation
strategy for recursive synthetic contamination.

%% --- For discussion.tex (new subsection or append to §Broader Impact) ---

\paragraph{Graduated versus binary filtering.}
The failure of graduated filtering to outperform binary
(Section~\ref{sec:graduated-filtering}) is not a null result awaiting
stronger evaluation---it reflects the information structure of recursive
synthetic text.
The mutual information $I(X_0; X_k)$ between original and
depth-$k$ text drops below an operational threshold at $\Kempirical$:
adjacent-depth discrimination beyond this boundary is weak (1v2 AUC
$\sim$0.61--0.70), and the residual ordinal signal that graduated
filtering attempts to exploit is insufficient to produce meaningful
downstream differentiation.
The graduated scheme's $\exp(-\alpha \cdot \hat{d})$ weighting assumes
a smooth quality gradient across depths, but the empirical feature
structure is step-like rather than smooth: quality metrics (surprisal
variance, JSD relative to human text) exhibit their largest discontinuity
at the human--synthetic boundary (depth~0$\to$1, JSD 0.21--0.29 dropping
to 0.005--0.016 for depth~1$\to$2), not at deeper transitions
(Figure~\ref{fig:feature_stepfunction}).
The graduated decay function is mismatched to this step structure---it
allocates weighting resolution across a gradient that does not exist.
This result also recontextualizes Drayson et al.'s binary SIR approach:
their use of binary filtering is not a methodological simplification
but already the optimal strategy, and our results provide the
information-theoretic explanation for why.
These conclusions are conditioned on LoRA fine-tuning (rank~8),
a 30{,}000-document corpus, and 0-shot evaluation on MMLU and
HellaSwag; full fine-tuning at larger scale with generation-intensive
tasks may probe finer curation differences.
